% Generated draft: requires mathematical and source review.
\documentclass[11pt]{article}
\usepackage[margin=25mm]{geometry}
\usepackage{fontspec}
\setmainfont{Latin Modern Roman}
\usepackage{amsmath,amssymb,mathtools,unicode-math}
\setmathfont{Latin Modern Math}
\usepackage{enumitem,needspace,xurl,hyperref}
\setlength{\parindent}{0pt}
\setlength{\parskip}{5pt}
\setlength{\emergencystretch}{3em}
\newcommand{\N}{\mathbb{N}}
\newcommand{\Z}{\mathbb{Z}}
\newcommand{\Q}{\mathbb{Q}}
\newcommand{\R}{\mathbb{R}}
\newcommand{\C}{\mathbb{C}}
\newcommand{\F}{\mathbb{F}}
\newcommand{\A}{\mathbb{A}}
\newcommand{\PP}{\mathbb{P}}
\newcommand{\E}{\mathbb{E}}
\newcommand{\eps}{\varepsilon}
\newcommand{\dd}{\,\mathrm{d}}
\begin{document}
\section*{UnsolvedMath ID: 30004054}
\textbf{Source problem code:} catalogue entry
\subsection{Definitions and mathematical statement}
This is a provenance-preserving statement-only draft for a composite record. The supplied registry marks its statement status as \texttt{unrecoverable}: the Latin-square clause is represented canonically by \texttt{OWR-16763-027}, and the bipartite-multigraph clause by \texttt{OWR-16763-030}. According to the supplied verification note, this extraction adds no distinct problem, and no separate clean statement is promoted. The two clauses below explicate the retained original wording; they are not asserted as proved results or as a newly recovered independent problem.

Let \(n\in\Z_{\geq 1}\), and write \([n]=\{1,\ldots,n\}\). A Latin square of order \(n\) is a function \(L:[n]\times[n]\to[n]\) such that each symbol occurs exactly once in each row and exactly once in each column. A partial transversal is a set \(T\subseteq[n]\times[n]\) such that, for any distinct \((i,j),(i',j')\in T\), one has \(i\ne i'\), \(j\ne j'\), and \(L(i,j)\ne L(i',j')\). A transversal is a partial transversal of cardinality \(n\).

\begin{enumerate}
\item \textbf{Retained Latin-square conjectures.} For every positive integer \(n\) and every Latin square \(L\) of order \(n\), there exists a partial transversal of cardinality \(n-1\). If \(n\) is odd, there exists a transversal of cardinality \(n\). For \(n=1\), the first assertion permits the empty partial transversal.
\item \textbf{Retained rainbow-matching conjecture.} Let \(G\) be a bipartite multigraph with vertex classes \(A\) and \(B\); parallel edges are allowed, and each edge joins a vertex of \(A\) to a vertex of \(B\). A proper edge-colouring with colour set \([n]\) is a function \(c:E(G)\to[n]\) for which distinct edges sharing a vertex have different colours. A matching is a set of edges no two of which share a vertex, and it is rainbow if its edges have pairwise distinct colours. For every positive integer \(n\), every such properly edge-coloured bipartite multigraph satisfying
\[
\lvert c^{-1}(i)\rvert\geq n+1\qquad\text{for every }i\in[n]
\]
has a matching \(M\subseteq E(G)\) satisfying
\[
\lvert M\cap c^{-1}(i)\rvert=1\qquad\text{for every }i\in[n].
\]
Thus \(M\) has \(n\) edges and uses every colour.
\end{enumerate}

The source does not explicitly specify finiteness of the multigraph. The displayed formulation imposes no additional finiteness restriction; its finite and unrestricted versions are equivalent, since one can retain exactly \(n+1\) edges of each of the finitely many colours before seeking the required matching. The supplied literature assessment concerns the second clause and does not provide a separate assessment of the Latin-square clauses.
\subsection{Short English statement}
The retained composite record asks whether every Latin square has a partial transversal missing only one row and, at odd order, a full transversal; it also asks whether \(n\) sufficiently large colour classes in a properly edge-coloured bipartite multigraph guarantee a matching using all \(n\) colours. The registry identifies both clauses as duplicates of other source records. Its supplied partially solved label and literature assessment require independent review.
\subsection{Sources}
\begin{itemize}
\item Local UnsolvedMath registry, \texttt{data/unsolvedmath/problems.json}, permanent ID \(30004054\), exact source problem code \texttt{catalogue entry}, title ``Composite Duplicate of Two Rainbow-Matching Conjectures.'' The supplied record retains the original composite statement, has no clean statement, and marks statement status as \texttt{unrecoverable}. Its verification note identifies canonical source records \texttt{OWR-16763-027} and \texttt{OWR-16763-030} and says that this composite extraction adds no distinct problem.
\item Supplied report citation: \emph{Graph Theory (2020)}, Oberwolfach Reports, DOI 10.4171/owr/2019/1, pp.~5--63. \url{https://doi.org/10.4171/owr/2019/1}. The registry gives proposed year \(2019\), but supplies no proposer. The background describes Oberwolfach Reports as documenting mathematical workshops and their open-problem discussions. The citation and provenance are supplied, not independently verified here.
\item Supplied literature source: Pokrovskiy, ``approximate Aharoni--Berger conjecture.'' \url{https://doi.org/10.1016/j.aim.2018.05.036}.
\item Supplied literature source: Keevash--Yepremyan, ``rainbow matchings.'' \url{https://arxiv.org/abs/1710.03041}.
\item Supplied literature assessment, checked exactly \texttt{2026-08-22}: the record is labelled \texttt{partially\_solved}. Its assessment says that the Aharoni--Berger rainbow-matching conjecture is not proved in full, credits Pokrovskiy with an \(n+o(n)\) asymptotic version, mentions related large-matching consequences, and describes the exact \(n+1\)-per-colour assertion as remaining open. This is supplied dated curation derived from the OWR material, not an independently verified literature assessment. The background explicitly requests independent verification of time-sensitive status and cited claims.
\end{itemize}
\end{document}
